\chapter*{Appendix}
\addcontentsline{toc}{chapter}{Appendix}
\label{chap:appendix}

\section*{A. Database Schema and Cypher Constraints}
The Neo4j property graph is fortified with unique entity constraints and semantic indexes to guarantee sub-millisecond AST-guided query execution and zero duplicate node anomalies.

\begin{lstlisting}[language=SQL, caption={Neo4j Database Uniqueness Constraints and Property Indexes}]
// Enforce unique identification across primary entities
CREATE CONSTRAINT unique_paper_id IF NOT EXISTS
FOR (p:Paper) REQUIRE p.id IS UNIQUE;

CREATE CONSTRAINT unique_dataset_name IF NOT EXISTS
FOR (d:Dataset) REQUIRE d.name IS UNIQUE;

CREATE CONSTRAINT unique_method_name IF NOT EXISTS
FOR (m:Method) REQUIRE m.name IS UNIQUE;

CREATE CONSTRAINT unique_metric_id IF NOT EXISTS
FOR (met:Metric) REQUIRE met.id IS UNIQUE;

CREATE CONSTRAINT unique_concept_name IF NOT EXISTS
FOR (c:Concept) REQUIRE c.name IS UNIQUE;

// Property and text indexes for rapid prefix and fulltext lookups
CREATE INDEX index_paper_title IF NOT EXISTS
FOR (p:Paper) ON (p.title);

CREATE INDEX index_method_category IF NOT EXISTS
FOR (m:Method) ON (m.category);

CREATE INDEX index_metric_value IF NOT EXISTS
FOR (met:Metric) ON (met.metric_name, met.value);
\end{lstlisting}

\section*{B. Extended Benchmark Questions and Ground Truth Verification}
Table~\ref{tab:extended_questions} details representative benchmark queries across all five cognitive tiers from our open-source 100-question academic dataset.

\begin{table}[htbp]
\centering
\small
\caption{Representative Benchmark Questions Across the 5 Cognitive Tiers}
\label{tab:extended_questions}
\begin{tabularx}{\textwidth}{l p{5.5cm} p{4.5cm} X}
\toprule
\textbf{ID} & \textbf{Question} & \textbf{Gold Target Reference} & \textbf{Cognitive Tier} \\
\midrule
\texttt{q01} & What is the BLEU score of Transformer (big) on WMT 2014 En-De? & 28.4 & Lookup \\
\texttt{q15} & What is the inference latency reduction of FlashAttention-2? & 2x speedup (up to 73\% FLOPs) & Lookup \\
\texttt{q21} & Which attention mechanism enabled BERT to perform bidirectional conditioning? & Multi-Head Self-Attention with masked tokens & Multi-Hop \\
\texttt{q38} & What mathematical operation in LoRA allows memory reduction during backward pass? & Low-rank matrix decomposition ($W_0 + \Delta W = W_0 + B \cdot A$) & Multi-Hop \\
\texttt{q41} & How many total attention heads does Transformer Base use across its encoder layers? & 48 heads (6 layers $\times$ 8 heads) & Aggregation \\
\texttt{q59} & What is the average parameter reduction percentage achieved by LoRA across transformer weights? & $> 99\%$ (rank $r=4$ or $r=8$) & Aggregation \\
\texttt{q61} & How does FlashAttention optimize GPU memory compared to the original Transformer? & IO-aware tiling avoiding HBM read/writes & Cross-Paper \\
\texttt{q77} & Contrast the pre-training objectives of BERT and RoBERTa. & RoBERTa removes NSP and uses dynamic masking & Cross-Paper \\
\texttt{q81} & Explain how FlashAttention redesigns attention computation to optimize GPU SRAM and HBM memory bandwidth. & Tiling algorithm splitting $Q, K, V$ into SRAM blocks with online softmax scaling & Synthesis \\
\texttt{q96} & Synthesize how LoRA achieves parameter-efficient fine-tuning without adding inference latency. & Low-rank matrix product $B \cdot A$ folded into weight $W$ during deployment ($W = W_0 + BA$) & Synthesis \\
\bottomrule
\end{tabularx}
\end{table}

\section*{C. Core Agent Cypher Generation Prompt Template}
The prompt below exemplifies the schema-constrained Cypher generation template deployed in the backend agent:

\begin{lstlisting}[language=Python, caption={Schema-Constrained Cypher Generation Prompt in Agent Pipeline}]
CYPHER_GENERATION_PROMPT = """
You are an expert Neo4j Cypher engineer querying an academic literature knowledge graph.
Given a user question, generate a clean, read-only Cypher query.

SCHEMA SPECIFICATION:
Nodes:
  - (:Paper {id, title, year, authors})
  - (:Method {id, name, category, description})
  - (:Metric {id, metric_name, value, unit, paper_id, split})
  - (:Dataset {id, name, modality, domain})
  - (:Concept {id, name, field})
  - (:Table {id, caption, paper_id, markdown_content})

Allowed Relationships:
  (:Paper)-[:USES_METHOD]->(:Method)
  (:Paper)-[:ACHIEVES]->(:Metric)
  (:Paper)-[:EVALUATED_ON]->(:Dataset)
  (:Paper)-[:DISCUSSES_CONCEPT]->(:Concept)
  (:Method)-[:ACHIEVES]->(:Metric)
  (:Metric)-[:ON_DATASET]->(:Dataset)

RULES:
1. ALWAYS use case-insensitive matching where appropriate: toLower(m.name) CONTAINS toLower('...')
2. DO NOT hallucinate edge types or node labels.
3. Return ONLY valid Cypher code enclosed in ```cypher ... ``` codeblocks.
4. Always include LIMIT 25 to prevent memory exhaustion.
Question: {question}
Cypher Query:
"""
\end{lstlisting}

\section*{D. System Environment and Configuration Parameters}
Table~\ref{tab:config_parameters} outlines the core configuration variables utilized across the production deployment.

\begin{table}[htbp]
\centering
\caption{Production System Configuration Parameters}
\label{tab:config_parameters}
\begin{tabular}{lll}
\toprule
\textbf{Configuration Variable} & \textbf{Value} & \textbf{Description} \\
\midrule
\texttt{LLM\_MODEL\_NAME} & \texttt{qwen2.5:7b-instruct-q4\_K\_M} & Primary generation model in Ollama \\
\texttt{EMBEDDING\_MODEL} & \texttt{nomic-embed-text} & 768-dim dense embedding generator \\
\texttt{NEO4J\_URI} & \texttt{bolt://localhost:7687} & Neo4j Bolt protocol endpoint \\
\texttt{NEO4J\_DATABASE} & \texttt{neo4j} & Primary graph database instance \\
\texttt{QDRANT\_HOST} & \texttt{localhost:6333} & Vector database REST / gRPC port \\
\texttt{CHUNK\_SIZE} & \texttt{512} & Dense passage chunk token size \\
\texttt{CHUNK\_OVERLAP} & \texttt{64} & Sliding window chunk overlap tokens \\
\texttt{VOTING\_ALPHA} & \texttt{0.65} & Weight allocated to graph verification \\
\texttt{API\_PORT} & \texttt{8000} & FastAPI server host port \\
\texttt{FRONTEND\_PORT} & \texttt{5173} & Vite React client development port \\
\bottomrule
\end{tabular}
\end{table}

\section*{E. GitHub Repository Information and Replication Guide}
The complete, end-to-end source code for the project, including backend APIs, frontend client interfaces, Neo4j ingestion scripts, and evaluation benchmarks, is hosted in our verified Git repository:

\begin{center}
\fbox{\parbox{0.9\textwidth}{\centering
\textbf{Project Repository:}\\
\url{https://github.com/THE-Adi-005/PaperEase_KG_FSD}\\
\vspace{0.2cm}
\small{\textit{Contains complete frontend, backend, graph seeds, 100-question benchmark dataset, and execution instructions.}}
}}
\end{center}

\subsection*{Replication Instructions}
To reproduce the experimental results and execute the full-stack system locally:
\begin{enumerate}
    \item \textbf{Clone the Repository:}
    \begin{lstlisting}[language=bash]
git clone https://github.com/THE-Adi-005/PaperEase_KG_FSD.git
cd PaperEase_KG_FSD
    \end{lstlisting}
    
    \item \textbf{Initialize Backend Services:}
    \begin{lstlisting}[language=bash]
cd backend
python -m venv venv
source venv/bin/activate  # On Windows: .\venv\Scripts\Activate.ps1
pip install -r requirements.txt
uvicorn app.main:app --host 0.0.0.0 --port 8000 --reload
    \end{lstlisting}
    
    \item \textbf{Initialize Neo4j and Ollama:}
    \begin{lstlisting}[language=bash]
# Ensure Neo4j is running on bolt://localhost:7687
ollama serve
ollama pull qwen2.5:7b-instruct-q4_K_M
ollama pull nomic-embed-text
    \end{lstlisting}
    
    \item \textbf{Initialize Frontend Interface:}
    \begin{lstlisting}[language=bash]
cd ../frontend
npm install
npm run dev
    \end{lstlisting}
    The interface will be accessible at \texttt{http://localhost:5173}.
    
    \item \textbf{Execute Benchmark Suite:}
    \begin{lstlisting}[language=bash]
cd ../backend
python ../eval/seed_experiment_results.py
    \end{lstlisting}
\end{enumerate}
